\documentclass[11pt]{article}
% Shared preamble for all daily exercises
% Update this file to change settings (padding, parindent, etc.) for all exercises

\usepackage[margin=0.75in]{geometry}
\usepackage{amsmath}
\usepackage{amssymb}

% Custom settings
\setlength{\parindent}{0pt}
\setlength{\parskip}{0.2cm}

\title{Daily Exercise}
\author{Dhruman Gupta}
\date{20th April}

\begin{document}

% Common header for daily exercises
% This file can be included in other LaTeX documents

\maketitle

\noindent
\textbf{Course:} CS-3330, Theory of Computation

\vspace{0.2cm}

\noindent
\textbf{Question:}\noindent 
For $k \geq 1$, define $\Sigma_k^P = \{ L \mid$ there exists a polynomial-time predicate $R$ such that 
\[
x \in L \iff \exists y_1 \forall y_2 \exists y_3 \cdots Q_k y_k. R(x, y_1, \dots, y_k)
\]
where the quantifiers alternate, starting with $\exists$, and each $y_i$ is polynomially bounded in $|x|$. Similarly, $\Pi_k^P$ starts with $\forall$.

Show that:
\begin{itemize}
    \item $\Sigma_1^P = NP$ and $\Pi_1^P = coNP$
    \item $\Sigma_k^P \subseteq \Sigma_{k+1}^P$ and $\Pi_k^P \subseteq \Pi_{k+1}^P$ for all $k$
    \item $\Sigma_k^P = co\Pi_k^P$
\end{itemize}

\textbf{Answer:}

First, for $k = 1$, the definition of $\Sigma_1^P$ is:

\[
x \in L \iff \exists y \text{ such that } |y| \leq p(|x|) \text{ and } R(x,y)
\]

for some polynomial $p$ and some polynomial-time predicate $R$.

But this is exactly the definition of $NP$: there exists a polynomially bounded certificate $y$ that can be verified in polynomial time. Hence,

\[
\Sigma_1^P = NP.
\]

Now, the definition of $\Pi_1^P$ is:

\[
x \in L \iff \forall y \text{ such that } |y| \leq p(|x|), R(x,y).
\]

Take complement. Then,

\[
x \notin L \iff \exists y \text{ such that } |y| \leq p(|x|) \text{ and } \neg R(x,y).
\]

Since $R$ is polynomial-time decidable, so is $\neg R$. Hence $\overline{L} \in NP$, and so $L \in coNP$. Therefore,

\[
\Pi_1^P = coNP.
\]

Now, let $L \in \Sigma_k^P$. Then there exists a polynomial-time predicate $R$ such that

\[
x \in L \iff \exists y_1 \forall y_2 \exists y_3 \cdots Q_k y_k \; R(x,y_1,\dots,y_k).
\]

To show that $L \in \Sigma_{k+1}^P$, just add one dummy variable at the end. Define

\[
R'(x,y_1,\dots,y_k,y_{k+1}) = R(x,y_1,\dots,y_k).
\]

Clearly $R'$ is still polynomial-time decidable. Then,

\[
x \in L \iff \exists y_1 \forall y_2 \exists y_3 \cdots Q_k y_k \, Q_{k+1} y_{k+1} \; R'(x,y_1,\dots,y_k,y_{k+1}).
\]

Since the extra variable does not affect the predicate, this is a valid $\Sigma_{k+1}^P$ definition. Hence,

\[
\Sigma_k^P \subseteq \Sigma_{k+1}^P.
\]

The proof that $\Pi_k^P \subseteq \Pi_{k+1}^P$ is identical: just add an extra dummy quantified variable at the end.

Finally, let $L \in \Sigma_k^P$. Then,

\[
x \in L \iff \exists y_1 \forall y_2 \exists y_3 \cdots Q_k y_k \; R(x,y_1,\dots,y_k).
\]

Take complement. Then,

\[
x \notin L \iff \neg(\exists y_1 \forall y_2 \exists y_3 \cdots Q_k y_k \; R(x,y_1,\dots,y_k)).
\]

Push the negation through the quantifiers. Every quantifier flips, so we get

\[
x \notin L \iff \forall y_1 \exists y_2 \forall y_3 \cdots \overline{Q_k} y_k \; \neg R(x,y_1,\dots,y_k).
\]

Since $R$ is polynomial-time decidable, so is $\neg R$. Hence $\overline{L} \in \Pi_k^P$, and so $L \in co\Pi_k^P$. Thus,

\[
\Sigma_k^P \subseteq co\Pi_k^P.
\]

Now let $L \in co\Pi_k^P$. Then $\overline{L} \in \Pi_k^P$, so

\[
x \in \overline{L} \iff \forall y_1 \exists y_2 \forall y_3 \cdots Q_k y_k \; R(x,y_1,\dots,y_k)
\]

for some polynomial-time predicate $R$.

Taking complement again,

\[
x \in L \iff \exists y_1 \forall y_2 \exists y_3 \cdots \overline{Q_k} y_k \; \neg R(x,y_1,\dots,y_k).
\]

Hence $L \in \Sigma_k^P$, and so

\[
co\Pi_k^P \subseteq \Sigma_k^P.
\]

Therefore,

\[
\Sigma_k^P = co\Pi_k^P.
\]

Similarly,

\[
\Pi_k^P = co\Sigma_k^P.
\]

Hence proved.

\end{document}
